\documentclass[11pt]{article}

% arXiv preprint style
\usepackage[utf8]{inputenc}
\usepackage[T1]{fontenc}
\usepackage{lmodern}
\usepackage[margin=1in]{geometry}
\usepackage{hyperref}
\usepackage{url}
\usepackage{booktabs}
\usepackage{amsfonts}
\usepackage{amsmath}
\usepackage{nicefrac}
\usepackage{microtype}
\usepackage{graphicx}
\usepackage{natbib}
\usepackage{doi}
\usepackage{xcolor}
\usepackage{caption}
\usepackage{subcaption}

\hypersetup{
    colorlinks=true,
    linkcolor=blue,
    citecolor=blue,
    urlcolor=blue
}

\title{The State of Blockchain Research (2013--2026)}

\author{
  Vedang Ratan Vatsa \\
  \texttt{vedangvats@gmail.com}
}

\date{May 2026}

\begin{document}

\maketitle

\begin{abstract}
This study presents a bibliometric analysis of blockchain and Web3 research from 2013 to 2026. A corpus of 100,024 academic papers, conference proceedings, whitepapers, and industry reports was compiled from Crossref, arXiv, Europe PMC, and OpenAlex. Using n-gram frequency analysis, year-over-year growth modeling, and citation distribution mapping, this study identifies research concentrations, shifting topic momentum, and the transition of the field from speculative interest to applied infrastructure engineering. The analysis yields three principal findings. First, the bigram ``blockchain technology'' appears 8,181 times, more than any other phrase, indicating that researchers treat blockchain as a component within larger systems rather than an isolated technology. Second, ``supply chain management'' is the most common trigram with 1,036 occurrences, showing that logistics and tracking applications exceed financial systems in publication volume. Third, the fastest-rising research topics between 2022 and 2026 focus on large language models, with ``language models'' growing 26x and ``large language'' growing 19x in blockchain-related papers. The dataset and source code are available for replication.
\end{abstract}

\noindent\textit{\textbf{Keywords:} blockchain, Web3, bibliometrics, n-gram analysis, research trends, supply chain, DeFi, tokenization}

\section{Introduction}
\label{sec:intro}

Blockchain research has expanded from early cryptographic cash concepts, such as Chaum's blind signatures \citep{chaum1983blind} and Dai's b-money \citep{dai1998bmoney}, to Nakamoto's peer-to-peer electronic cash system \citep{nakamoto2008bitcoin}, and now to tens of thousands of publications annually. However, volume alone does not clarify the direction of the field. A researcher entering the discipline must understand where the academic community is focusing its efforts and how those priorities are changing.

This paper addresses this question through an empirical study of 100,024 documents spanning 2013 to 2026. The dataset was compiled from four major academic indexes and curated industry reports. We perform n-gram frequency analysis on document titles, measure year-over-year publication growth, map citation distributions, and track specific topic cohorts to observe how research focus has shifted.

Rather than predicting future breakthroughs, this work seeks to describe the historical and current focus of the research community with precision.

\section{Related Work}
\label{sec:related}

Previous bibliometric studies of blockchain research have analyzed smaller datasets. Firdaus et al. \citep{firdaus2019root} evaluated 3,580 papers from Scopus, while Miau and Yang \citep{miau2018bibliometrics} examined early Bitcoin research trends across 1,119 publications. Casino et al. \citep{casino2019systematic} surveyed applications across 314 papers. Early synthetic literature by Swan \citep{swan2015blockchain} and Pilkington \citep{pilkington2016blockchain} categorized the conceptual transition from digital currency to decentralized application platforms. Zheng et al. \citep{zheng2017overview} provided an architectural overview of consensus mechanisms, which build upon classic distributed systems theories like the Byzantine Generals Problem \citep{lamport1982byzantine} and Practical Byzantine Fault Tolerance \citep{castro1999practical}. The present study expands on this work by evaluating a corpus of 100,024 documents, covering three additional years of rapid growth (2024 to 2026) and applying n-gram trend analysis to identify shifting thematic cohorts.

\section{Dataset and Methodology}
\label{sec:method}

\subsection{Data Collection}

Documents were collected from four major scholarly indexes, using API services including the Crossref REST API \citep{crossref2026api} and the OpenAlex database \citep{openalex2022open}, as shown in Table~\ref{tab:sources}.

\begin{table}[h]
\centering
\caption{Data sources and document counts.}
\label{tab:sources}
\begin{tabular}{lrl}
\toprule
Source & Documents & Type \\
\midrule
Crossref API & 88,037 & Journal articles \\
Europe PMC & 6,624 & Conference papers \\
arXiv & 1,949 & Preprints \\
Curated institutional & 3,414 & Reports, whitepapers \\
\midrule
\textbf{Total} & \textbf{100,024} & \\
\bottomrule
\end{tabular}
\end{table}

We queried the sources using a whitelist of 52 blockchain-related keywords (including blockchain, cryptocurrency, DeFi, NFT, smart contract, and tokenization). The results were filtered against a list of 62 terms from unrelated fields, such as cardiology, astrophysics, and photovoltaics, to remove false positives. A manual review of a random sample of 50 documents yielded a 100\% relevance rate.

\subsection{Analysis Methods}

\textbf{N-gram extraction.} Document titles were lowercased, stripped of HTML formatting, and tokenized. We applied a stopword list of 53 common English words to remove noise. Unigrams (single words), bigrams (two-word phrases), and trigrams (three-word phrases) were then extracted and counted across the corpus.

\textbf{Year-over-year growth.} Documents were grouped by publication year to calculate annual growth rates as the percentage change relative to the preceding year.

\textbf{Citation analysis.} Citation counts were extracted from the Crossref database. We calculated the median, mean, and 99th percentile citation thresholds.

\textbf{Trend detection.} Bigram frequencies were compared between the 2022 and 2026 publication cohorts. Topics with at least 10 occurrences in 2026 and a growth ratio exceeding 5x were classified as rising topics.

\section{Results}
\label{sec:results}

\subsection{Publication Volume and Growth}

Blockchain research output has grown annually since 2015, as shown in Table~\ref{tab:yearly}.

\begin{table}[h]
\centering
\caption{Yearly publication volume and citation statistics.}
\label{tab:yearly}
\begin{tabular}{lrrl}
\toprule
Year & Documents & YoY Growth & Avg Citations \\
\midrule
2015 & 788 & -- & 64.1 \\
2016 & 897 & +13.8\% & 107.6 \\
2017 & 1,584 & +76.6\% & 91.2 \\
2018 & 3,752 & +136.9\% & 70.8 \\
2019 & 6,122 & +63.2\% & 44.1 \\
2020 & 7,008 & +14.5\% & 21.2 \\
2021 & 8,993 & +28.3\% & 13.9 \\
2022 & 11,263 & +25.2\% & 8.2 \\
2023 & 14,462 & +28.4\% & 5.4 \\
2024 & 16,394 & +13.4\% & 3.9 \\
2025 & 20,781 & +26.8\% & 1.0 \\
2026* & 7,980 & -- & 0.1 \\
\bottomrule
\end{tabular}
\end{table}

\noindent *2026 data is partial (January to May).

Two periods of rapid acceleration are evident. The first occurred in 2017 and 2018, when annual output rose by 136.9\%, matching the growth of the initial ICO phase and Bitcoin's run to \$20,000. The second period, spanning 2023 to 2025, showed steady annual growth between 13\% and 28\%. This latter expansion was driven by institutional developments, including BlackRock's BUIDL fund, spot Bitcoin ETFs, MiCA regulations \citep{eu2023mica}, and real-world asset tokenization platforms \citep{rwa2026dashboard}, rather than speculative retail activity.

The inverse relationship between publication volume and average citations is expected, as older papers have had more time to accumulate references. However, the 2016 cohort remains the highest-impact group, with an average of 107.6 citations, containing early, highly cited smart contract and Ethereum architectural studies \citep{buterin2014ethereum, wood2014ethereum}.

\subsection{Keyword Frequency}

The ten most frequent single keywords across the 100,024 document titles are shown in Table~\ref{tab:keywords}.

\begin{table}[h]
\centering
\caption{Top 10 single keywords by frequency.}
\label{tab:keywords}
\begin{tabular}{clrr}
\toprule
Rank & Keyword & Count & \% of docs \\
\midrule
1 & blockchain & 54,780 & 54.8\% \\
2 & decentralized & 12,897 & 12.9\% \\
3 & technology & 9,982 & 10.0\% \\
4 & smart & 8,043 & 8.0\% \\
5 & data & 8,032 & 8.0\% \\
6 & bitcoin & 6,024 & 6.0\% \\
7 & security & 5,702 & 5.7\% \\
8 & IoT & 5,521 & 5.5\% \\
9 & cryptocurrency & 5,116 & 5.1\% \\
10 & supply & 4,783 & 4.8\% \\
\bottomrule
\end{tabular}
\end{table}

While ``blockchain'' appears in 54.8\% of all titles, ``bitcoin'' is present in only 6.0\%, indicating that the research community has expanded well beyond the original cryptocurrency framework. The prominence of ``IoT'' at rank 8 (5,521 mentions) reflects a substantial body of research focusing on blockchain-secured device networks \citep{christidis2016blockchains, fernandez2018review, reyna2018blockchain}.

\subsection{Bigram and Trigram Analysis}

The five most frequent bigrams and trigrams are shown in Table~\ref{tab:ngrams}.

\begin{table}[h]
\centering
\caption{Top 5 bigrams and trigrams by frequency.}
\label{tab:ngrams}
\begin{tabular}{clr|clr}
\toprule
\# & Bigram & Count & \# & Trigram & Count \\
\midrule
1 & blockchain technology & 8,181 & 1 & supply chain management & 1,036 \\
2 & supply chain & 3,960 & 2 & central bank digital & 959 \\
3 & smart contract & 2,887 & 3 & bank digital currency & 852 \\
4 & blockchain enabled & 2,313 & 4 & distributed ledger technology & 643 \\
5 & smart contracts & 2,216 & 5 & non fungible tokens & 558 \\
\bottomrule
\end{tabular}
\end{table}

``Supply chain management'' is the most frequent trigram in the corpus, with 1,036 occurrences. This makes logistics and traceability the largest individual application area by publication volume. ``Central bank digital'' (959) and ``bank digital currency'' (852) represent a combined block of nearly 1,800 papers focused on central bank digital currencies (CBDCs), establishing state-backed digital currencies as the second-largest research theme.

\subsection{Blockchain-ML Convergence}

The intersection of blockchain and machine learning represents a substantial portion of the corpus. ``Federated learning'' appears 1,712 times \citep{yang2019federated, mcmahan2017communication}, while ``machine learning'' (1,779) and ``privacy preserving'' (1,342) \citep{kosba2016hawk} define a cluster of over 4,800 papers exploring the integration of blockchain with machine learning models.

\begin{table}[h]
\centering
\caption{Blockchain-ML convergence trigrams.}
\label{tab:blml}
\begin{tabular}{lr}
\toprule
Trigram & Count \\
\midrule
blockchain federated learning & 382 \\
blockchain machine learning & 306 \\
federated learning blockchain & 300 \\
\bottomrule
\end{tabular}
\end{table}

These papers describe systems where blockchain operates as a coordination layer for distributed machine learning, allowing multiple parties to train shared models on local, private data without a central intermediary \citep{lu2020blockchain}. This convergence was rarely observed in publications before 2021.

\subsection{Rising Topics (2022 vs 2026)}

Table~\ref{tab:rising} shows the fastest-growing bigrams between the 2022 and 2026 publication cohorts.

\begin{table}[h]
\centering
\caption{Fastest-rising bigrams from 2022 to 2026.}
\label{tab:rising}
\begin{tabular}{lrr}
\toprule
Bigram & 2026 count & Growth vs 2022 \\
\midrule
language models & 26 & 26.0x \\
enhancing transparency & 25 & 25.0x \\
large language & 19 & 19.0x \\
chain transparency & 18 & 18.0x \\
quantum resistant & 15 & 15.0x \\
secure transparent & 36 & 12.0x \\
\bottomrule
\end{tabular}
\end{table}

The rising topics fall into three primary areas: the application of large language models in blockchain systems, an emphasis on transaction and supply chain transparency driven by regulatory frameworks such as MiCA \citep{eu2023mica} and the GENIUS Act \citep{us2025genius}, and early work on quantum-resistant cryptography \citep{merkle1987digital} to secure protocols against future decryption threats.

\subsection{Citation Distribution}

The citation distribution is heavily right-skewed. Among the 100,024 documents, 51.8\% have at least one citation. The median citation count is 1, the mean is 14.7, and the top 1\% threshold is 249 citations. The most-cited document is Nakamoto's Bitcoin whitepaper \citep{nakamoto2008bitcoin}, which has 14,286 citations.

The top 20 keywords account for 22.8\% of all word occurrences across 36,199 unique words, suggesting a field that has diversified into specialized domains while retaining a core focus on smart contracts, infrastructure security, and logistics.

\section{Discussion}
\label{sec:discussion}

\subsection{From Speculation to Infrastructure}

The data reveals a clear evolution. Between 2015 and 2018, blockchain research was dominated by cryptocurrency-focused topics, including mining, consensus mechanisms, and market speculation, matching the 2017 to 2018 growth spike.

Since 2020, the thematic composition has shifted. ``Supply chain'' and ``IoT'' now outrank ``cryptocurrency'' in keyword frequency. The emergence of terms like ``smart contract vulnerability'' (288 trigram occurrences) and ``access control'' (1,180 bigram occurrences) indicates that researchers are focused on hardening existing protocols for production use rather than exploring theoretical concepts.

\subsection{The Federated Learning Bridge}

The overlap between blockchain and machine learning represents a major structural shift in the literature. This combination is mathematically practical: federated learning requires a coordination mechanism for distributed model training, and blockchain provides an auditable, trustless ledger for that coordination.

Rather than copying sensitive records into a single central database, which poses severe privacy risks, federated learning leaves the data secure on local devices and only shares model adjustments, using the blockchain as a secure, neutral coordinator. This architecture has immediate applications in healthcare, where patient records must remain within hospital networks, financial services, where institutions train shared fraud detection models without exposing customer accounts, and logistics optimization \citep{saberi2019blockchain, tian2016agri}.

\subsection{The CBDC Research Wave}

Central bank digital currencies represent a distinct research cluster of approximately 1,800 papers. This work is driven by public sector pilots, including the European Central Bank's digital euro project \citep{ecb2024euro}, the People's Bank of China's e-CNY trials, and the Bank of England's consultations on a digital pound.

The volume of CBDC research reflects an institutional commitment. These papers focus on solving specific engineering challenges, including offline transaction capability, privacy-preserving monitoring, and interoperability protocols between national systems \citep{auer2021cbdcs}.

\subsection{Limitations}

Title-level bibliometric analysis has limitations. It cannot evaluate the internal quality or methodology of individual papers. A paper titled ``Blockchain for Supply Chain Management'' might represent a rigorous empirical study or a superficial review. While citation counts provide a proxy for impact, they are a lagging indicator that naturally favors older publications.

Additionally, citation metadata from Crossref can be incomplete as some publishers do not report citation indexes, and self-citations are not isolated. Finally, the 2026 cohort represents only partial-year data through May.

\section{Conclusion}
\label{sec:conclusion}

An analysis of 100,024 documents reveals a research field that has transitioned beyond speculative cryptocurrency applications. Academic output in 2026 is dominated by infrastructure development, including supply chain traceability, IoT device management, protocol security, and central bank digital currencies. The convergence with federated learning represents an active research frontier. The rising emphasis on transparency, showing up to 25x growth, and quantum resistance, showing 15x growth, indicates a field preparing for regulatory compliance and long-term cryptographic durability.

The complete dataset is available at \url{https://veda.ng/web3-reports}.

\section*{References}

\begin{thebibliography}{32}
\providecommand{\natexlab}[1]{#1}
\providecommand{\url}[1]{\texttt{#1}}
\expandafter\ifx\csname urlstyle\endcsname\relax
  \providecommand{\doi}[1]{doi: #1}\else
  \providecommand{\doi}{doi: \begingroup \urlstyle{rm}\Url}\fi

\bibitem[Nakamoto(2008)]{nakamoto2008bitcoin}
S.~Nakamoto.
\newblock Bitcoin: A peer-to-peer electronic cash system.
\newblock \emph{Whitepaper}, 2008.

\bibitem[Buterin(2014)]{buterin2014ethereum}
V.~Buterin.
\newblock Ethereum: A next-generation smart contract and decentralized
  application platform.
\newblock \emph{Whitepaper}, 2014.

\bibitem[Christidis and Devetsikiotis(2016)]{christidis2016blockchains}
K.~Christidis and M.~Devetsikiotis.
\newblock Blockchains and smart contracts for the internet of things.
\newblock \emph{IEEE Access}, 4:\penalty0 2292--2303, 2016.

\bibitem[Swan(2015)]{swan2015blockchain}
M.~Swan.
\newblock \emph{Blockchain: Blueprint for a New Economy}.
\newblock O'Reilly Media, 2015.

\bibitem[Szabo(1997)]{szabo1997formalizing}
N.~Szabo.
\newblock Formalizing and securing relationships on public networks.
\newblock \emph{First Monday}, 2\penalty0 (9), 1997.

\bibitem[Lamport et~al.(1982)Lamport, Shostak, and Pease]{lamport1982byzantine}
L.~Lamport, R.~Shostak, and M.~Pease.
\newblock The Byzantine generals problem.
\newblock \emph{ACM Transactions on Programming Languages and Systems},
  4\penalty0 (3):\penalty0 382--401, 1982.

\bibitem[Castro and Liskov(1999)]{castro1999practical}
M.~Castro and B.~Liskov.
\newblock Practical Byzantine fault tolerance.
\newblock In \emph{Proceedings of OSDI}, pages 173--186, 1999.

\bibitem[Wood(2014)]{wood2014ethereum}
G.~Wood.
\newblock Ethereum: A secure decentralised generalised transaction ledger.
\newblock \emph{Ethereum Project Yellow Paper}, 2014.

\bibitem[Luu et~al.(2016)Luu, Chu, Olickel, Saxena, and Hobor]{luu2016making}
L.~Luu, D.-H. Chu, H.~Olickel, P.~Saxena, and A.~Hobor.
\newblock Making smart contracts smarter.
\newblock In \emph{Proceedings of ACM CCS}, pages 254--269, 2016.

\bibitem[Kosba et~al.(2016)Kosba, Miller, Shi, Wen, and
  Papamanthou]{kosba2016hawk}
A.~Kosba, A.~Miller, E.~Shi, Z.~Wen, and C.~Papamanthou.
\newblock Hawk: The blockchain model of cryptography and privacy-preserving
  smart contracts.
\newblock In \emph{Proceedings of IEEE S\&P}, pages 839--858, 2016.

\bibitem[Saberi et~al.(2019)Saberi, Kouhizadeh, Sarkis, and
  Shen]{saberi2019blockchain}
S.~Saberi, M.~Kouhizadeh, J.~Sarkis, and L.~Shen.
\newblock Blockchain technology and its relationships to sustainable supply
  chain management.
\newblock \emph{International Journal of Production Research}, 57\penalty0
  (7):\penalty0 2117--2135, 2019.

\bibitem[Tian(2016)]{tian2016agri}
F.~Tian.
\newblock An agri-food supply chain traceability system for china based on rfid
  and blockchain technology.
\newblock In \emph{Proceedings of International Conference on Service Systems
  and Service Management}, pages 1--6, 2016.

\bibitem[Fernandez-Carames and Fraga-Lamas(2018)]{fernandez2018review}
T.~M. Fernandez-Carames and P.~Fraga-Lamas.
\newblock A review on the use of blockchain for the internet of things.
\newblock \emph{IEEE Access}, 6:\penalty0 32979--33001, 2018.

\bibitem[Reyna et~al.(2018)Reyna, Martin, Chen, Soler, and
  Diaz]{reyna2018blockchain}
A.~Reyna, C.~Martin, J.~Chen, E.~Soler, and M.~Diaz.
\newblock On blockchain and its integration with iot: Challenges and
  opportunities.
\newblock \emph{Future Generation Computer Systems}, 88:\penalty0 173--190,
  2018.

\bibitem[Yang et~al.(2019)Yang, Liu, Chen, and Tong]{yang2019federated}
Q.~Yang, Y.~Liu, T.~Chen, and Y.~Tong.
\newblock Federated machine learning: Concept and applications.
\newblock \emph{ACM Transactions on Intelligent Systems and Technology},
  10\penalty0 (2):\penalty0 1--19, 2019.

\bibitem[McMahan et~al.(2017)]{mcmahan2017communication}
B.~McMahan et~al.
\newblock Communication-efficient learning of deep networks from decentralized
  data.
\newblock In \emph{Proceedings of AISTATS}, pages 1273--1282, 2017.

\bibitem[Lu et~al.(2020)Lu, Huang, Zhang, Maharjan, and Zhang]{lu2020blockchain}
Y.~Lu, X.~Huang, K.~Zhang, S.~Maharjan, and Y.~Zhang.
\newblock Blockchain empowered asynchronous federated learning for secure data
  sharing in internet of vehicles.
\newblock \emph{IEEE Transactions on Vehicular Technology}, 69\penalty0
  (4):\penalty0 4298--4311, 2020.

\bibitem[Dai(1998)]{dai1998bmoney}
W.~Dai.
\newblock b-money.
\newblock \emph{Proposal}, 1998.

\bibitem[Chaum(1983)]{chaum1983blind}
D.~Chaum.
\newblock Blind signatures for untraceable payments.
\newblock In \emph{Advances in Cryptology}, pages 199--203, 1983.

\bibitem[Merkle(1987)]{merkle1987digital}
R.~C. Merkle.
\newblock A digital signature based on a conventional encryption function.
\newblock In \emph{Proceedings of CRYPTO}, pages 369--378, 1987.

\bibitem[Auer(2021)]{auer2021cbdcs}
M.~Auer.
\newblock Cbdcs beyond borders: Results from a survey of central banks.
\newblock \emph{Bank for International Settlements, BIS Papers No. 116}, 2021.

\bibitem[European~Central~Bank(2024)]{ecb2024euro}
European~Central~Bank.
\newblock Digital euro project.
\newblock \emph{Report}, 2024.

\bibitem[European~Union(2023)]{eu2023mica}
European~Union.
\newblock Regulation (eu) 2023/1114 (markets in crypto-assets, mica).
\newblock \emph{Official Journal of the European Union}, 2023.

\bibitem[U.S.~Congress(2025)]{us2025genius}
U.S.~Congress.
\newblock Genius act of 2025.
\newblock \emph{S.394, 119th Congress}, 2025.

\bibitem[Firdaus et~al.(2019)Firdaus, Razak, Feizollah, Hashem, Hazim, and
  Imran]{firdaus2019root}
A.~Firdaus, N.~Razak, M.~Feizollah, I.~Hashem, M.~Hazim, and M.~Imran.
\newblock The root causes of the blockchain research: A bibliometric analysis.
\newblock \emph{Telematics and Informatics}, 36:\penalty0 46--54, 2019.

\bibitem[Miau and Yang(2018)]{miau2018bibliometrics}
S.~Miau and J.~M. Yang.
\newblock Bibliometrics-based evaluation of the blockchain research trend.
\newblock \emph{Technology Analysis and Strategic Management}, 30\penalty0
  (9):\penalty0 1029--1045, 2018.

\bibitem[Casino et~al.(2019)Casino, Dasaklis, and
  Patsakis]{casino2019systematic}
F.~Casino, T.~K. Dasaklis, and C.~Patsakis.
\newblock A systematic literature review of blockchain-based applications:
  Current status, classification and open issues.
\newblock \emph{Telematics and Informatics}, 36:\penalty0 55--81, 2019.

\bibitem[Zheng et~al.(2017)Zheng, Xie, Dai, Chen, and Wang]{zheng2017overview}
Z.~Zheng, S.~Xie, H.~Dai, X.~Chen, and H.~Wang.
\newblock An overview of blockchain technology: Architecture, consensus, and
  future trends.
\newblock In \emph{Proceedings of IEEE BigData Congress}, pages 557--564, 2017.

\bibitem[Pilkington(2016)]{pilkington2016blockchain}
M.~Pilkington.
\newblock Blockchain technology: Principles and applications.
\newblock In \emph{Research Handbook on Digital Transformations}. Edward Elgar
  Publishing, 2016.

\bibitem[RWA.xyz(2026)]{rwa2026dashboard}
RWA.xyz.
\newblock Tokenized asset dashboard.
\newblock \emph{Market Report}, 2026.
\newblock \url{https://www.rwa.xyz}.

\bibitem[Crossref(2026)]{crossref2026api}
Crossref.
\newblock Crossref rest api.
\newblock \emph{Technical Documentation}, 2026.
\newblock \url{https://api.crossref.org}.

\bibitem[OpenAlex(2022)]{openalex2022open}
OpenAlex.
\newblock Openalex: A fully-open index of scholarly works.
\newblock \emph{Database Documentation}, 2022.
\newblock \url{https://openalex.org}.

\end{thebibliography}

\end{document}
